\section{CSS3视觉效果}\label{css3ux89c6ux89c9ux6548ux679c}

\subsection{1. 圆角与边框}\label{ux5706ux89d2ux4e0eux8fb9ux6846}

CSS3引入了多种边框特性，使设计师无需使用图片即可创建丰富的边框效果，这些特性在智慧水利平台界面设计中能显著提升视觉吸引力。

\subsubsection{1.1 border-radius属性}\label{border-radiusux5c5eux6027}

\passthrough{\lstinline!border-radius!}属性允许设计师创建圆角元素，从而使界面更加现代化、友好。

\paragraph{基本用法}\label{ux57faux672cux7528ux6cd5}

\begin{lstlisting}
/* 所有角相同的圆角 */
.card {
    border-radius: 8px;
    background-color: white;
    padding: 20px;
    box-shadow: 0 2px 8px rgba(0, 0, 0, 0.1);
}

/* 指定不同角的圆角 */
.panel-top {
    border-radius: 8px 8px 0 0; /* 上左、上右、下右、下左 */
    background-color: #1890ff;
    color: white;
    padding: 10px 15px;
}

/* 椭圆形圆角 */
.pill-button {
    border-radius: 50px;
    padding: 8px 20px;
    background-color: #1890ff;
    color: white;
    border: none;
}

/* 创建圆形 */
.avatar {
    border-radius: 50%;
    width: 60px;
    height: 60px;
    object-fit: cover;
}
\end{lstlisting}

\paragraph{在智慧水利平台中的应用}\label{ux5728ux667aux6167ux6c34ux5229ux5e73ux53f0ux4e2dux7684ux5e94ux7528}

\begin{lstlisting}
/* 数据卡片圆角 */
.data-card {
    border-radius: 12px;
    background-color: white;
    padding: 20px;
    margin-bottom: 20px;
    box-shadow: 0 2px 10px rgba(0, 0, 0, 0.08);
}

/* 状态指示器圆角 */
.status-badge {
    display: inline-block;
    padding: 4px 12px;
    border-radius: 20px;
    font-size: 12px;
    font-weight: 500;
}

.status-normal {
    background-color: #f6ffed;
    color: #52c41a;
}

.status-warning {
    background-color: #fffbe6;
    color: #faad14;
}

.status-danger {
    background-color: #fff2f0;
    color: #ff4d4f;
}
\end{lstlisting}

\subsubsection{1.2 box-shadow属性}\label{box-shadowux5c5eux6027}

\passthrough{\lstinline!box-shadow!}属性允许为元素添加阴影效果，增加界面的深度感和层次感。

\paragraph{基本用法}\label{ux57faux672cux7528ux6cd5-1}

\begin{lstlisting}
/* 基础阴影 */
.card {
    box-shadow: 0 2px 8px rgba(0, 0, 0, 0.1);
}

/* 多重阴影 */
.floating-panel {
    box-shadow: 0 2px 4px rgba(0, 0, 0, 0.1),
                0 8px 16px rgba(0, 0, 0, 0.1);
}

/* 内阴影 */
.inset-panel {
    box-shadow: inset 0 2px 4px rgba(0, 0, 0, 0.1);
}

/* 有色阴影 */
.primary-button {
    box-shadow: 0 4px 12px rgba(24, 144, 255, 0.4);
}
\end{lstlisting}

\paragraph{在智慧水利平台中的应用}\label{ux5728ux667aux6167ux6c34ux5229ux5e73ux53f0ux4e2dux7684ux5e94ux7528-1}

\begin{lstlisting}
/* 水平面板悬浮效果 */
.dashboard-card {
    background-color: white;
    border-radius: 12px;
    padding: 20px;
    box-shadow: 0 2px 10px rgba(0, 0, 0, 0.08);
    transition: transform 0.3s ease, box-shadow 0.3s ease;
}

.dashboard-card:hover {
    transform: translateY(-5px);
    box-shadow: 0 10px 20px rgba(0, 0, 0, 0.12);
}

/* 按钮阴影 */
.primary-button {
    background-color: #1890ff;
    color: white;
    border: none;
    padding: 10px 20px;
    border-radius: 4px;
    box-shadow: 0 2px 6px rgba(24, 144, 255, 0.4);
    transition: background-color 0.3s, box-shadow 0.3s;
}

.primary-button:hover {
    background-color: #40a9ff;
    box-shadow: 0 4px 12px rgba(24, 144, 255, 0.6);
}

/* 警告卡片阴影 */
.alert-card {
    border-left: 4px solid #ff4d4f;
    background-color: #fff2f0;
    padding: 15px;
    margin-bottom: 15px;
    border-radius: 0 4px 4px 0;
    box-shadow: 0 2px 8px rgba(255, 77, 79, 0.2);
}
\end{lstlisting}

\subsubsection{1.3 border-image属性}\label{border-imageux5c5eux6027}

\passthrough{\lstinline!border-image!}属性允许使用图像作为元素的边框，创建复杂的边框效果。

\paragraph{基本用法}\label{ux57faux672cux7528ux6cd5-2}

\begin{lstlisting}
.decorated-panel {
    border: 15px solid transparent;
    border-image: url('border-pattern.png') 30 round;
    padding: 20px;
}
\end{lstlisting}

\paragraph{在智慧水利平台中的应用}\label{ux5728ux667aux6167ux6c34ux5229ux5e73ux53f0ux4e2dux7684ux5e94ux7528-2}

\begin{lstlisting}
/* 水纹边框 */
.water-card {
    border: 10px solid transparent;
    border-image: url('/images/water-border.png') 20 repeat;
    padding: 20px;
    background-color: #f0f9ff;
}

/* 渐变边框 */
.gradient-border-card {
    border: 5px solid;
    border-image: linear-gradient(45deg, #1890ff, #69c0ff) 1;
    padding: 20px;
}
\end{lstlisting}

\subsection{2. 背景与渐变}\label{ux80ccux666fux4e0eux6e10ux53d8}

CSS3大幅增强了背景处理能力，包括多重背景、背景尺寸控制和渐变背景等特性，使界面设计更加灵活和富有表现力。

\subsubsection{2.1 多重背景}\label{ux591aux91cdux80ccux666f}

CSS3允许为元素设置多个背景图像，这些背景图像按照声明顺序从上到下层叠。

\begin{lstlisting}
.header-section {
    background-image: 
        url('logo-watermark.png'),
        linear-gradient(rgba(0, 0, 0, 0.5), rgba(0, 0, 0, 0.5)),
        url('header-bg.jpg');
    background-position: 
        center right 20px,
        center,
        center;
    background-repeat: 
        no-repeat,
        repeat,
        no-repeat;
    background-size: 
        auto 60%,
        cover,
        cover;
    height: 300px;
    color: white;
    padding: 40px;
}
\end{lstlisting}

\subsubsection{2.2
背景尺寸与定位}\label{ux80ccux666fux5c3aux5bf8ux4e0eux5b9aux4f4d}

CSS3引入了更精确控制背景尺寸和位置的新属性。

\begin{lstlisting}
/* 背景尺寸 */
.cover-bg {
    background-image: url('map.jpg');
    background-size: cover; /* 覆盖整个容器 */
    background-position: center;
}

.contain-bg {
    background-image: url('icon.png');
    background-size: contain; /* 保持比例缩放以适应容器 */
    background-repeat: no-repeat;
    background-position: center;
}

/* 具体尺寸 */
.fixed-bg {
    background-image: url('pattern.png');
    background-size: 50px 50px;
    background-repeat: repeat;
}

/* 背景定位 */
.positioned-bg {
    background-image: url('marker.png');
    background-size: 30px 30px;
    background-repeat: no-repeat;
    background-position: right 20px bottom 10px; /* 距右边20px，距底部10px */
}
\end{lstlisting}

\subsubsection{2.3 线性渐变}\label{ux7ebfux6027ux6e10ux53d8}

线性渐变创建沿着直线方向颜色平滑过渡的图像。

\begin{lstlisting}
/* 基础线性渐变 - 从上到下 */
.vertical-gradient {
    background: linear-gradient(#1890ff, #0050b3);
}

/* 指定方向 */
.angle-gradient {
    background: linear-gradient(45deg, #1890ff, #0050b3);
}

/* 多色渐变 */
.multi-color-gradient {
    background: linear-gradient(to right, #1890ff, #69c0ff, #bae7ff);
}

/* 指定渐变位置 */
.positioned-gradient {
    background: linear-gradient(to bottom, 
        #1890ff 0%, 
        #69c0ff 40%, 
        #bae7ff 100%);
}

/* 硬边界渐变 */
.hard-stop-gradient {
    background: linear-gradient(to right,
        #1890ff 50%,
        #0050b3 50%);
}
\end{lstlisting}

\subsubsection{2.4 径向渐变}\label{ux5f84ux5411ux6e10ux53d8}

径向渐变创建从中心点辐射的颜色过渡效果。

\begin{lstlisting}
/* 基础径向渐变 */
.radial-gradient {
    background: radial-gradient(#1890ff, #0050b3);
}

/* 指定形状和大小 */
.ellipse-gradient {
    background: radial-gradient(ellipse at center, #1890ff, #0050b3);
}

.circle-gradient {
    background: radial-gradient(circle at center, #1890ff, #0050b3);
}

/* 指定位置 */
.positioned-radial {
    background: radial-gradient(circle at top right, #1890ff, #0050b3);
}

/* 多色径向渐变 */
.multi-color-radial {
    background: radial-gradient(circle,
        #bae7ff 0%,
        #69c0ff 50%,
        #1890ff 100%);
}
\end{lstlisting}

\subsubsection{2.5 重复渐变}\label{ux91cdux590dux6e10ux53d8}

重复渐变可以创建条纹和纹理效果。

\begin{lstlisting}
/* 重复线性渐变 - 创建条纹 */
.striped-background {
    background: repeating-linear-gradient(
        45deg,
        #e6f7ff,
        #e6f7ff 10px,
        #bae7ff 10px,
        #bae7ff 20px
    );
}

/* 重复径向渐变 - 创建水波纹效果 */
.ripple-background {
    background: repeating-radial-gradient(
        circle at center,
        #e6f7ff,
        #e6f7ff 10px,
        #bae7ff 10px,
        #bae7ff 20px
    );
}
\end{lstlisting}

\subsubsection{2.6
在智慧水利平台中的应用}\label{ux5728ux667aux6167ux6c34ux5229ux5e73ux53f0ux4e2dux7684ux5e94ux7528-3}

\begin{lstlisting}
/* 水位监测面板背景 */
.water-level-panel {
    background: linear-gradient(to bottom, #e6f7ff, #ffffff);
    border-radius: 12px;
    padding: 20px;
    box-shadow: 0 2px 10px rgba(0, 0, 0, 0.08);
}

/* 流量指示器背景 */
.flow-indicator {
    height: 60px;
    border-radius: 30px;
    background: repeating-linear-gradient(
        90deg,
        #bae7ff,
        #bae7ff 10px,
        #e6f7ff 10px,
        #e6f7ff 20px
    );
    background-size: 200% 100%;
    animation: flowAnimation 10s linear infinite;
    position: relative;
}

@keyframes flowAnimation {
    0% { background-position: 0% 0%; }
    100% { background-position: 200% 0%; }
}

/* 降雨强度图 */
.rainfall-intensity-map {
    height: 300px;
    width: 300px;
    border-radius: 50%;
    background: radial-gradient(
        circle,
        #ffffff 0%,
        #bae7ff 40%,
        #69c0ff 70%,
        #1890ff 85%,
        #0050b3 100%
    );
    box-shadow: 0 4px 20px rgba(24, 144, 255, 0.3);
}

/* 仪表盘标题背景 */
.dashboard-header {
    background-image: 
        url('/images/logo-watermark.png'),
        linear-gradient(90deg, #1890ff, #096dd9);
    background-position: 
        right 20px center,
        center;
    background-repeat: 
        no-repeat,
        no-repeat;
    background-size: 
        auto 60%,
        cover;
    color: white;
    padding: 20px;
    border-radius: 8px 8px 0 0;
}
\end{lstlisting}

\subsection{3. 文本效果}\label{ux6587ux672cux6548ux679c}

CSS3引入了多种增强文本表现力的特性，使文字内容更具可读性和视觉吸引力。

\subsubsection{3.1 文本阴影}\label{ux6587ux672cux9634ux5f71}

\passthrough{\lstinline!text-shadow!}属性为文本添加阴影效果，增强文字的立体感和可读性。

\begin{lstlisting}
/* 基础文本阴影 */
.title {
    text-shadow: 2px 2px 4px rgba(0, 0, 0, 0.3);
}

/* 发光文本效果 */
.glow-text {
    color: white;
    text-shadow: 0 0 10px rgba(24, 144, 255, 0.8);
}

/* 多重阴影 */
.multi-shadow-text {
    color: #1890ff;
    text-shadow: 1px 1px 2px rgba(0, 0, 0, 0.3),
                 0 0 8px rgba(24, 144, 255, 0.5);
}

/* 文字浮雕效果 */
.embossed-text {
    color: #444;
    text-shadow: 1px 1px 1px white;
}

.pressed-text {
    color: #444;
    text-shadow: -1px -1px 1px white;
}
\end{lstlisting}

\subsubsection{3.2 文本溢出}\label{ux6587ux672cux6ea2ux51fa}

CSS3提供了更好的方法来处理文本溢出情况。

\begin{lstlisting}
/* 单行文本溢出显示省略号 */
.truncate {
    white-space: nowrap;
    overflow: hidden;
    text-overflow: ellipsis;
    max-width: 200px;
}

/* 多行文本溢出 */
.multi-line-truncate {
    display: -webkit-box;
    -webkit-line-clamp: 3; /* 显示行数 */
    -webkit-box-orient: vertical;
    overflow: hidden;
    max-height: 4.5em; /* 行高 × 行数 */
}
\end{lstlisting}

\subsubsection{3.3 字体特性}\label{ux5b57ux4f53ux7279ux6027}

CSS3增强了对字体的控制能力，包括网页字体和字体特性控制。

\begin{lstlisting}
/* 使用@font-face导入自定义字体 */
@font-face {
    font-family: 'WaterIcons';
    src: url('/fonts/water-icons.woff2') format('woff2'),
         url('/fonts/water-icons.woff') format('woff');
    font-weight: normal;
    font-style: normal;
}

.water-icon {
    font-family: 'WaterIcons';
}

/* 字体平滑 */
body {
    -webkit-font-smoothing: antialiased;
    -moz-osx-font-smoothing: grayscale;
}

/* 字体特性控制 */
.tabular-nums {
    font-variant-numeric: tabular-nums; /* 等宽数字，适合数据表格 */
}
\end{lstlisting}

\subsubsection{3.4 文本修饰}\label{ux6587ux672cux4feeux9970}

CSS3提供了更多文本修饰选项。

\begin{lstlisting}
/* 文本下划线 */
.custom-underline {
    text-decoration: underline;
    text-decoration-color: #1890ff;
    text-decoration-thickness: 2px;
    text-decoration-style: wavy;
}

/* 文本描边 */
.outlined-text {
    -webkit-text-stroke: 1px black;
    color: white;
}
\end{lstlisting}

\subsubsection{3.5
在智慧水利平台中的应用}\label{ux5728ux667aux6167ux6c34ux5229ux5e73ux53f0ux4e2dux7684ux5e94ux7528-4}

\begin{lstlisting}
/* 数据面板标题 */
.panel-title {
    font-size: 20px;
    font-weight: 500;
    color: #1890ff;
    text-shadow: 0 1px 2px rgba(24, 144, 255, 0.2);
    margin-bottom: 16px;
}

/* 站点名称溢出处理 */
.station-name {
    font-weight: 500;
    white-space: nowrap;
    overflow: hidden;
    text-overflow: ellipsis;
    max-width: 200px;
}

/* 警告文本 */
.warning-text {
    color: #faad14;
    font-weight: 500;
    text-shadow: 0 0 4px rgba(250, 173, 20, 0.4);
}

.danger-text {
    color: #ff4d4f;
    font-weight: 600;
    text-shadow: 0 0 4px rgba(255, 77, 79, 0.4);
}

/* 数据值显示 */
.data-value {
    font-family: 'DIN Condensed', sans-serif;
    font-variant-numeric: tabular-nums;
    font-size: 24px;
    font-weight: 500;
    color: #1890ff;
}

/* 降雨量大值特殊显示 */
.rainfall-heavy {
    color: #096dd9;
    font-weight: bold;
    text-decoration: underline;
    text-decoration-color: #096dd9;
    text-decoration-thickness: 2px;
}
\end{lstlisting}

\subsection{4. 颜色与透明度}\label{ux989cux8272ux4e0eux900fux660eux5ea6}

CSS3提供了多种颜色表示方法，包括RGBA、HSLA以及透明度控制。

\subsubsection{4.1 RGBA颜色}\label{rgbaux989cux8272}

RGBA允许设置颜色的同时控制其透明度。

\begin{lstlisting}
/* RGBA颜色 - RGB + Alpha通道 */
.overlay {
    background-color: rgba(0, 0, 0, 0.5); /* 半透明黑色 */
}

.highlight {
    background-color: rgba(24, 144, 255, 0.2); /* 淡蓝色高亮 */
    border: 1px solid rgba(24, 144, 255, 0.5);
}
\end{lstlisting}

\subsubsection{4.2 HSLA颜色}\label{hslaux989cux8272}

HSLA基于色相、饱和度、亮度和透明度定义颜色，更加直观。

\begin{lstlisting}
/* HSLA颜色 - 色相、饱和度、亮度、透明度 */
.primary-button {
    background-color: hsla(210, 100%, 50%, 1); /* 蓝色 */
}

.primary-button:hover {
    background-color: hsla(210, 100%, 55%, 1); /* 稍亮的蓝色 */
}

.overlay-panel {
    background-color: hsla(0, 0%, 0%, 0.7); /* 70%不透明黑色 */
}
\end{lstlisting}

\subsubsection{4.3 透明度控制}\label{ux900fux660eux5ea6ux63a7ux5236}

CSS3允许通过多种方式控制元素透明度。

\begin{lstlisting}
/* 使用opacity属性 - 影响整个元素及其子元素 */
.faded {
    opacity: 0.7;
}

/* 使用rgba背景 - 只影响背景色 */
.translucent-bg {
    background-color: rgba(255, 255, 255, 0.8);
}

/* 使用透明色 */
.transparent-border {
    border: 1px solid transparent;
}
\end{lstlisting}

\subsubsection{4.4
在智慧水利平台中的应用}\label{ux5728ux667aux6167ux6c34ux5229ux5e73ux53f0ux4e2dux7684ux5e94ux7528-5}

\begin{lstlisting}
/* 数据图表提示框 */
.chart-tooltip {
    background-color: rgba(0, 0, 0, 0.75);
    color: white;
    padding: 10px;
    border-radius: 4px;
    box-shadow: 0 2px 8px rgba(0, 0, 0, 0.2);
}

/* 水位预警指示器 */
.water-level-indicator {
    position: relative;
}

.water-level-indicator::after {
    content: "";
    position: absolute;
    top: 0;
    left: 0;
    right: 0;
    bottom: 0;
    background-color: rgba(255, 77, 79, 0.2);
    border-radius: inherit;
    animation: pulse 2s infinite;
}

@keyframes pulse {
    0% { opacity: 0.2; }
    50% { opacity: 0.5; }
    100% { opacity: 0.2; }
}

/* 叠加信息面板 */
.map-overlay-panel {
    background-color: rgba(255, 255, 255, 0.9);
    backdrop-filter: blur(5px); /* 模糊背景效果 */
    border-radius: 8px;
    padding: 15px;
    box-shadow: 0 2px 10px rgba(0, 0, 0, 0.1);
}

/* 状态色彩系统 */
:root {
    --normal-color: hsla(135, 60%, 39%, 1);
    --normal-bg: hsla(135, 60%, 95%, 1);
    --warning-color: hsla(38, 100%, 50%, 1);
    --warning-bg: hsla(38, 100%, 95%, 1);
    --danger-color: hsla(0, 100%, 65%, 1);
    --danger-bg: hsla(0, 100%, 96%, 1);
}

.status-normal {
    color: var(--normal-color);
    background-color: var(--normal-bg);
}

.status-warning {
    color: var(--warning-color);
    background-color: var(--warning-bg);
}

.status-danger {
    color: var(--danger-color);
    background-color: var(--danger-bg);
}
\end{lstlisting}

\subsection{5. 高级视觉效果}\label{ux9ad8ux7ea7ux89c6ux89c9ux6548ux679c}

CSS3提供了多种创建高级视觉效果的功能，如模糊、滤镜和混合模式等。

\subsubsection{5.1 滤镜效果}\label{ux6ee4ux955cux6548ux679c}

CSS3的\passthrough{\lstinline!filter!}属性允许应用图形滤镜效果，如模糊、亮度调整、对比度等。

\begin{lstlisting}
/* 模糊效果 */
.blur {
    filter: blur(5px);
}

/* 灰度效果 */
.grayscale {
    filter: grayscale(100%);
}

/* 亮度调整 */
.brighten {
    filter: brightness(150%);
}

/* 对比度调整 */
.increase-contrast {
    filter: contrast(180%);
}

/* 阴影效果 */
.drop-shadow {
    filter: drop-shadow(3px 3px 5px rgba(0, 0, 0, 0.3));
}

/* 多重滤镜组合 */
.vintage {
    filter: sepia(50%) contrast(120%) brightness(90%);
}
\end{lstlisting}

\subsubsection{5.2
背景模糊效果}\label{ux80ccux666fux6a21ux7ccaux6548ux679c}

使用\passthrough{\lstinline!backdrop-filter!}属性可以对元素背后的内容应用模糊或其他滤镜效果。

\begin{lstlisting}
/* 背景模糊面板 */
.glass-panel {
    background-color: rgba(255, 255, 255, 0.2);
    backdrop-filter: blur(10px);
    padding: 20px;
    border-radius: 10px;
    border: 1px solid rgba(255, 255, 255, 0.3);
}
\end{lstlisting}

\subsubsection{5.3 混合模式}\label{ux6df7ux5408ux6a21ux5f0f}

CSS3的混合模式允许指定元素内容或背景如何与底层元素混合。

\begin{lstlisting}
/* 混合模式 - 将内容与背景混合 */
.multiply-blend {
    background-image: url('texture.jpg');
    background-color: #1890ff;
    background-blend-mode: multiply;
}

/* 混合模式 - 将元素与其下方元素混合 */
.overlay-blend {
    mix-blend-mode: overlay;
}
\end{lstlisting}

\subsubsection{5.4 渐变蒙版}\label{ux6e10ux53d8ux8499ux7248}

使用渐变和混合模式创建蒙版效果。

\begin{lstlisting}
/* 渐变蒙版 */
.fade-mask {
    position: relative;
}

.fade-mask::after {
    content: "";
    position: absolute;
    bottom: 0;
    left: 0;
    right: 0;
    height: 100px;
    background: linear-gradient(transparent, white);
    pointer-events: none; /* 允许点击穿透 */
}
\end{lstlisting}

\subsubsection{5.5
在智慧水利平台中的应用}\label{ux5728ux667aux6167ux6c34ux5229ux5e73ux53f0ux4e2dux7684ux5e94ux7528-6}

\begin{lstlisting}
/* 水库照片灰度效果（历史数据） */
.historical-image {
    filter: grayscale(80%) contrast(110%);
    transition: filter 0.3s ease;
}

.historical-image:hover {
    filter: grayscale(0%) contrast(100%);
}

/* 天气状态指示器 */
.weather-icon.rainy {
    filter: drop-shadow(0 0 2px rgba(0, 110, 255, 0.5));
}

.weather-icon.sunny {
    filter: brightness(120%) saturate(120%);
}

/* 模糊状态面板 */
.glass-status-panel {
    background-color: rgba(255, 255, 255, 0.15);
    backdrop-filter: blur(10px);
    border-radius: 12px;
    padding: 20px;
    border: 1px solid rgba(255, 255, 255, 0.2);
    color: white;
}

/* 卫星图像增强 */
.satellite-image {
    filter: contrast(120%) saturate(110%) brightness(105%);
}

/* 热力图效果 */
.heatmap-overlay {
    background-image: url('/images/heatmap-gradient.png');
    mix-blend-mode: multiply;
    opacity: 0.8;
}

/* 图表数据区域渐变蒙版 */
.chart-container {
    position: relative;
}

.chart-scroll-indicator {
    position: absolute;
    right: 0;
    top: 0;
    bottom: 0;
    width: 50px;
    background: linear-gradient(to right, transparent, white);
    pointer-events: none;
}
\end{lstlisting}

通过这些CSS3视觉效果，智慧水利平台的界面可以变得更加现代化、专业化，并且能够更有效地传达信息和提升用户体验。圆角、阴影和渐变等特性不仅仅是装饰，它们也能引导用户注意力、表达关系层级，并使复杂信息更加直观易懂。
